\section*{第 \ref{ch:g10:freshwater-resources} 章——淡水资源}

\begin{solution}{exo:g10:freshwater-resources:1}
\omterm{def:g10:freshwater-resources:stock}{储量}是指某一时刻储库中所含的水量，单位为
\unit{km^3}（或千克）。\omterm{def:g10:freshwater-resources:stock}{通量}是指单位时间内由一个
储库转移到另一个储库的水量，单位为 \unit{km^3} 每年。\omterm{def:g10:freshwater-resources:residence-time}{停留时间}
是水在某一储库中平均停留的时间，单位为天或年；在
稳态下，它等于\omterm{def:g10:freshwater-resources:stock}{储量}除以\omterm{def:g10:freshwater-resources:stock}{通量}。
\end{solution}

\begin{solution}{exo:g10:freshwater-resources:2}
约 \qty{2.5}{\percent} 的水是淡水。淡水中
\qty{68.7}{\percent} 以冰盖、冰川和永久积雪的形式冻结，另有
\qty{30.1}{\percent} 储存于地下。河流中的水量最少，
仅 \num{2120}~\unit{km^3}。
\end{solution}

\begin{solution}{exo:g10:freshwater-resources:3}
\omterm{def:g10:freshwater-resources:renewable}{可再生淡水}是每年由降水补给更新的淡水，
并以地表径流和地下水的形式流经某一地区：是按年计的水量，
即\omterm{def:g10:freshwater-resources:stock}{通量}。\omterm{def:g10:freshwater-resources:fossil-groundwater}{化石地下水}是很久以前进入含水层的水，
而今天的补给无法使其更新：它是只能被逐步抽取的\omterm{def:g10:freshwater-resources:stock}{储量}，
即\omterm{def:g10:freshwater-resources:stock}{储量}。
\end{solution}

\begin{solution}{exo:g10:freshwater-resources:4}
$91\,000/2120 \approx 43$ 倍。
\end{solution}

\begin{solution}{exo:g10:freshwater-resources:5}
大气中的水更新最快（约 \num{9.5} 天），海洋
在所计算的储库中更新最慢（约 \num{3200} 年）；深层
地下水更为古老。$3200 \times 365.25/9.5 \approx \num{120000}$
倍之久。
\end{solution}

\begin{solution}{exo:g10:freshwater-resources:6}
流入：$413\,000 + 73\,000 = \num{486000}$~\unit{km^3} 每年，因此
$T = 12\,700/486\,000 \approx \num{0.0261}$ 年 $\approx \num{9.5}$ 天。
流出：$373\,000 + 113\,000 = \num{486000}$~\unit{km^3}，两者相同：大气处于
稳态，\omterm{def:g10:freshwater-resources:residence-time}{停留时间}相同。
\end{solution}

\begin{solution}{exo:g10:freshwater-resources:7}
$1\,335\,040/413 \approx \num{3230}$ 年。稳态下，海洋
输入与输出相等：$373 + 40 = 413$ 千~\unit{km^3} 每年，
因此由输入量算得的\omterm{def:g10:freshwater-resources:residence-time}{停留时间}相同。
\end{solution}

\begin{solution}{exo:g10:freshwater-resources:8}
陆地每年接收 \num{113000}~\unit{km^3}；其中 \num{73000} 经
蒸发和蒸腾作用返回大气，\num{40000} 以径流形式流入海洋：
$73\,000 + 40\,000 = \num{113000}$。返回大气的比例为
$73/113 \approx \qty{65}{\percent}$。
\end{solution}

\begin{solution}{exo:g10:freshwater-resources:9}
$273.2 \times 10^6 \times 1233.5 \approx \qty{3.37e11}{m^3}$，即
\qty{337}{km^3}：几乎相当于咸海在
1987 年仍保有的全部水量。
\end{solution}

\begin{solution}{exo:g10:freshwater-resources:10}
$15.8 \times 0.3048 \approx \qty{4.82}{m}$。在 1950 年至
2015 年的 65 年间，$482/65 \approx \qty{7.4}{cm}$ 每年。
\end{solution}

\begin{solution}{exo:g10:freshwater-resources:11}
$2/0.01 = 200$ 年。\omterm{def:g10:freshwater-resources:stock}{储量}现在每年减少 $0.05 - 0.01 =
\qty{0.04}{km^3}$，因此可维持 $2/0.04 = 50$ 年。
\end{solution}

\begin{solution}{exo:g10:freshwater-resources:12}
1960 年：$1093/71.1 \approx 15$ 年；2010 年：$98.1/11.4 \approx 8.6$ 年。
蒸发每年从湖面每平方千米带走大致相同厚度的水层；
湖泊缩小时，面积随之缩小，
蒸发的总水量也随之减少。
\end{solution}

\begin{solution}{exo:g10:freshwater-resources:13}
(1) 净损失为 $1.2 - 0.2 = \qty{1.0}{m}$ 每年，作用于
面积为 $100 \times 10^6$~\unit{m^2} 的湖面：$10^8$~\unit{m^3}，\qty{0.1}{km^3} 每年必须
流入。(2) 流入量降至 \qty{0.05}{km^3}，湖泊损失
$\qty{0.05}{km^3} = 5 \times 10^7$~\unit{m^3}，相当于一层
$5 \times 10^7/10^8 = \qty{0.5}{m}$。(3) 水位下降时，岸线
露出水面，湖泊面积缩小：水面减少，蒸发也随之减少，
直至蒸发与流入量重新达到平衡。
\end{solution}

\begin{solution}{exo:g10:freshwater-resources:14}
含水层储水 $100 \times 0.20 = \qty{20}{m}$，即
\qty{20000}{mm}：按每年 \qty{1}{mm} 计，需 \num{20000} 年。撒哈拉含水层的深层水
年龄为 \num{200000} 至一百万年，约为前者的十至五十倍：
它是在较湿润的过去气候期间获得补给的。
\end{solution}

\begin{solution}{exo:g10:freshwater-resources:15}
$337/453\,000 \approx \qty{0.00074}{km}$，约合水层 \qty{0.74}{m}。水位
下降了 \qty{4.8}{m}，约为其六倍半。水
只充填岩石的孔隙（并非全部都能排出）：当
地下水位下降一米时，释放出的水层远薄于
一米，此处约为 $0.74/4.8 \approx \num{0.15}$~m。
\end{solution}

\begin{solution}{pb:g10:freshwater-resources:1}
\textbf{1.} $1093/68\,000 \approx \qty{0.016}{km}$，约 \qty{16}{m}。

\textbf{2.} $1093 \times 10^9 \approx \qty{1.09e12}{m^3}$。

\textbf{3.} \qty{10}{g/L} 即 \qty{10}{kg/m^3}：$1.093 \times 10^{12} \times 10
\approx \qty{1.1e13}{kg}$，约 \num{11} 十亿吨。

\textbf{4.} $1093/71.1 \approx 15$ 年。

\textbf{5.} 没有河流流出它：其水只能以蒸气形式离开（或渗入
地下），而盐分留了下来。

\textbf{6.} $71.1/68\,000 \approx \qty{0.00105}{km}$，约 \qty{1.05}{m}。

\textbf{7.} $9.4/68\,000 \approx \qty{0.00014}{km}$，约 \qty{0.14}{m}。

\textbf{8.} $71.1 - 9.4 = \qty{61.7}{km^3}$ 每年。

\textbf{9.} $79 + 37 = \qty{116}{km^3}$：$61.7/116 \approx \qty{53}{\percent}$，
约一半。

\textbf{10.} 输入量 $1.0 + 6.2 = \qty{7.2}{km^3}$；输出量 \qty{36.8}{km^3}；
净变化 $7.2 - 36.8 = \qty{-29.6}{km^3}$。

\textbf{11.} $(502.7 - 327.2)/5 = \qty{35.1}{km^3}$ 每年，与
1987 年收支中的 \qty{29.6}{km^3} 量级相同（个别年份有所不同：1984 年蒸发量为
\qty{47.9}{km^3}）。

\textbf{12.} $36.8/0.00105 \approx \qty{35000}{km^2}$。

\textbf{13.} $35\,000/68\,000 \approx \qty{52}{\percent}$，约一半。

\textbf{14.} 湖泊已经萎缩：
湖面面积减半，蒸发的水量也约减半。

\textbf{15.} $345.6/1093 \approx \qty{32}{\percent}$，约三分之一。

\textbf{16.} $10 \times 1093/345.6 \approx \qty{31.6}{g/L}$。

\textbf{17.} 实测盐度偏低：部分盐分已离开
水体，在出露的湖床上结晶，并被风以含盐尘埃的形式
带走。

\textbf{18.} 每年蒸发的厚度约为 \qty{1.05}{m}。

\textbf{19.} $7.2/35\,000 \approx \qty{0.0002}{km}$，约 \qty{0.20}{m}。

\textbf{20.} $1.05 - 0.20 \approx 0.85$（未取整时为 \num{0.84}）：
\textbf{1987 年咸海水位下降约 \qty{0.8}{m}，
接近一年内下降一米}。
\end{solution}
